\documentclass[12pt,a4paper]{article}

\usepackage[utf8]{inputenc}
\usepackage[vietnamese]{babel}
\usepackage{amsmath,amssymb,amsthm}
\usepackage{geometry}
\usepackage{enumitem}
\usepackage[most]{tcolorbox}
\usepackage{hyperref}
\usepackage{fancyhdr}
\usepackage{tikz}
\usepackage{eso-pic}
\usepackage{xcolor}

\geometry{a4paper, margin=2cm, bottom=2.5cm}
\hypersetup{colorlinks=true, linkcolor=blue!70!black}
\setlist[itemize]{topsep=4pt, itemsep=2pt}
\setlist[enumerate]{topsep=4pt, itemsep=2pt}

\AddToShipoutPictureFG{%
  \AtPageCenter{%
    \begin{tikzpicture}[remember picture, overlay]
      \node[rotate=45, scale=1, opacity=0.06]{\fontsize{60}{72}\selectfont\bfseries\sffamily Đặng Tấn Quí};
    \end{tikzpicture}%
  }%
}

\pagestyle{fancy}
\fancyhf{}
\fancyhead[L]{\small\textit{Toán kinh tế 1}}
\fancyhead[R]{\small\textit{Đặng Tấn Quí}}
\fancyfoot[C]{\thepage}
\fancyfoot[L]{\footnotesize\textcopyright\ Đặng Tấn Quí}
\fancyfoot[R]{\footnotesize SĐT: 0377\,122\,210}
\renewcommand{\headrulewidth}{0.4pt}
\renewcommand{\footrulewidth}{0.4pt}

\fancypagestyle{plain}{\fancyhf{}\fancyfoot[C]{\thepage}\fancyfoot[L]{\footnotesize\textcopyright\ Đặng Tấn Quí}\fancyfoot[R]{\footnotesize SĐT: 0377\,122\,210}\renewcommand{\headrulewidth}{0pt}\renewcommand{\footrulewidth}{0.4pt}}

\newtcolorbox{dinhnghia}[1][]{enhanced, breakable, colback=blue!3!white, colframe=blue!60!black, fonttitle=\bfseries, title={Định nghĩa}, #1}
\newtcolorbox{dinhly}[1][]{enhanced, breakable, colback=green!3!white, colframe=green!50!black, fonttitle=\bfseries, title={Định lý}, #1}
\newtcolorbox{tinhchat}[1][]{enhanced, breakable, colback=yellow!5!white, colframe=orange!50!black, fonttitle=\bfseries, title={Tính chất}, #1}
\newtcolorbox{phuongphap}[1][]{enhanced, breakable, colback=gray!5!white, colframe=gray!50!black, fonttitle=\bfseries, title={Phương pháp}, #1}
\newtcolorbox{luuy}[1][]{enhanced, breakable, colback=red!3!white, colframe=red!50!black, fonttitle=\bfseries, title={Lưu ý}, #1}
\newtcolorbox{congthuc}[1][]{enhanced, breakable, colback=cyan!3!white, colframe=cyan!50!black, fonttitle=\bfseries, title={Công thức}, #1}
\newtcolorbox{vidu}[1][]{enhanced, breakable, colback=violet!3!white, colframe=violet!50!black, fonttitle=\bfseries, title={Ví dụ}, #1}

\begin{document}

\thispagestyle{empty}
\begin{tikzpicture}[remember picture, overlay]
  \fill[blue!80!black] (current page.south west) rectangle (current page.north east);
  \node[anchor=west, white, font=\fontsize{26}{32}\bfseries\selectfont]
    at ([xshift=3cm, yshift=-7.5cm]current page.north west) {TOÁN KINH TẾ 1};
  \node[anchor=west, white, font=\fontsize{18}{24}\selectfont]
    at ([xshift=3cm, yshift=-9cm]current page.north west) {Hàm KT, đạo hàm, Lagrange, Leontief, lãi kép};
  \node[anchor=west, white, font=\Large\bfseries] at ([xshift=3cm, yshift=-13cm]current page.north west) {Biên soạn:};
  \node[anchor=west, yellow!80!white, font=\fontsize{22}{28}\bfseries\selectfont] at ([xshift=3cm, yshift=-14.5cm]current page.north west) {ĐẶNG TẤN QUÍ};
  \node[anchor=south east, white, font=\Large\bfseries] at ([xshift=-2cm, yshift=3cm]current page.south east) {\the\year};
\end{tikzpicture}

\newpage
\tableofcontents
\newpage

\section{Hàm số trong kinh tế}

\subsection{Các hàm kinh tế cơ bản}

\begin{dinhnghia}
  \textbf{Hàm cầu} $Q_D = f(P)$: lượng cầu phụ thuộc giá. Thường giảm khi $P$ tăng.
\end{dinhnghia}

\begin{dinhnghia}
  \textbf{Hàm cung} $Q_S = g(P)$: lượng cung phụ thuộc giá. Thường tăng khi $P$ tăng.
\end{dinhnghia}

\begin{dinhnghia}
  \textbf{Hàm doanh thu} $R = R(Q) = P \cdot Q$: doanh thu theo sản lượng.
\end{dinhnghia}

\begin{dinhnghia}
  \textbf{Hàm chi phí} $C = C(Q)$: tổng chi phí sản xuất. \textbf{Hàm chi phí bình quân}: $AC = \dfrac{C(Q)}{Q}$.
\end{dinhnghia}

\begin{dinhnghia}
  \textbf{Hàm lợi nhuận} $\pi = R(Q) - C(Q)$.
\end{dinhnghia}

\subsection{Điểm cân bằng thị trường}

\begin{dinhnghia}
  \textbf{Điểm cân bằng} $(P^*, Q^*)$: $Q_D(P^*) = Q_S(P^*)$.
\end{dinhnghia}

\section{Đạo hàm và ứng dụng}

\subsection{Đạo hàm cơ bản}

\begin{congthuc}
  $(x^n)' = nx^{n-1}$, $(e^x)' = e^x$, $(\ln x)' = \dfrac{1}{x}$, $(uv)' = u'v + uv'$, $\left(\dfrac{u}{v}\right)' = \dfrac{u'v - uv'}{v^2}$.
\end{congthuc}

\subsection{Ý nghĩa kinh tế của đạo hàm}

\begin{dinhnghia}
  \textbf{Chi phí biên} $MC = C'(Q)$: chi phí tăng thêm khi sản xuất thêm 1 đơn vị.
\end{dinhnghia}

\begin{dinhnghia}
  \textbf{Doanh thu biên} $MR = R'(Q)$: doanh thu tăng thêm khi bán thêm 1 đơn vị.
\end{dinhnghia}

\begin{dinhnghia}
  \textbf{Hệ số co giãn} của $Q$ theo $P$: $E = \dfrac{P}{Q} \cdot \dfrac{dQ}{dP} = \dfrac{dQ/Q}{dP/P}$.
\end{dinhnghia}

\begin{luuy}
  $|E| > 1$: cầu co giãn; $|E| < 1$: cầu ít co giãn; $|E| = 1$: cầu co giãn đơn vị.
\end{luuy}

\subsection{Tối ưu hóa}

\begin{phuongphap}[title={Tìm mức sản lượng tối đa hóa lợi nhuận}]
  \begin{enumerate}
    \item $\pi(Q) = R(Q) - C(Q)$.
    \item $\pi'(Q) = 0$ $\Rightarrow$ $MR = MC$ (điều kiện cần).
    \item Kiểm tra $\pi''(Q) < 0$ (điều kiện đủ).
  \end{enumerate}
\end{phuongphap}

\begin{dinhly}
  Lợi nhuận đạt cực đại khi doanh thu biên bằng chi phí biên: $MR = MC$.
\end{dinhly}

\section{Hàm nhiều biến trong kinh tế}

\begin{dinhnghia}[title={Đạo hàm riêng và ý nghĩa}]
  $f(x_1, \ldots, x_n)$: sản lượng phụ thuộc nhiều yếu tố.

  $\dfrac{\partial f}{\partial x_i}$: năng suất biên của yếu tố $x_i$ (giữ các yếu tố khác cố định).
\end{dinhnghia}

\begin{dinhnghia}[title={Hàm sản xuất Cobb--Douglas}]
  $Q = AK^\alpha L^\beta$, $A > 0$, $\alpha, \beta > 0$.
  \begin{itemize}
    \item $\alpha + \beta = 1$: hiệu suất không đổi theo quy mô.
    \item $\alpha + \beta > 1$: hiệu suất tăng; $< 1$: giảm.
  \end{itemize}
\end{dinhnghia}

\begin{dinhnghia}[title={Hàm tiện ích}]
  $U(x_1, x_2)$: mức thỏa dụng. Đường bàng quan: $U(x_1, x_2) = c$. Tỷ lệ thay thế biên:
  \[
    MRS_{12} = -\frac{dx_2}{dx_1}\bigg|_{U = c} = \frac{\partial U/\partial x_1}{\partial U/\partial x_2}.
  \]
\end{dinhnghia}

\begin{phuongphap}[title={Tối ưu có ràng buộc --- Lagrange}]
  $\max U(x_1, x_2)$ t.c. $p_1 x_1 + p_2 x_2 = M$.

  $\mathcal{L} = U(x_1, x_2) - \lambda(p_1 x_1 + p_2 x_2 - M)$.

  Điều kiện: $\dfrac{\partial U/\partial x_1}{p_1} = \dfrac{\partial U/\partial x_2}{p_2} = \lambda$ (tiện ích biên trên giá bằng nhau).
\end{phuongphap}

\section{Ma trận và mô hình kinh tế}

\begin{dinhnghia}[title={Hệ phương trình tuyến tính trong kinh tế}]
  $\mathbf{A}\mathbf{x} = \mathbf{b}$: mô hình cân bằng thị trường nhiều hàng hóa, hệ IS-LM tuyến tính.
\end{dinhnghia}

\begin{dinhnghia}[title={Mô hình đầu vào--đầu ra (Leontief)}]
  $\mathbf{x} = \mathbf{A}\mathbf{x} + \mathbf{d}$ $\Rightarrow$ $\mathbf{x} = (\mathbf{I} - \mathbf{A})^{-1}\mathbf{d}$,

  $\mathbf{A}$: ma trận hệ số kỹ thuật ($a_{ij}$: lượng ngành $i$ cần cho 1 đơn vị ngành $j$), $\mathbf{d}$: cầu cuối cùng.
\end{dinhnghia}

\begin{dinhly}[title={Điều kiện Hawkins--Simon}]
  $(\mathbf{I} - \mathbf{A})^{-1} \ge \mathbf{0}$ (nghiệm không âm) khi và chỉ khi mọi định thức con chính của $\mathbf{I} - \mathbf{A}$ đều dương.
\end{dinhly}

\section{Chuỗi thời gian rời rạc cơ bản}

\begin{dinhnghia}[title={Dãy số kinh tế}]
  Dãy $\{y_t\}$: GDP, CPI, giá cổ phiếu theo kỳ. Tốc độ tăng trưởng: $g_t = \dfrac{y_t - y_{t-1}}{y_{t-1}}$.
\end{dinhnghia}

\begin{congthuc}[title={Lãi kép rời rạc}]
  $FV = PV(1 + r)^n$, $PV = \dfrac{FV}{(1 + r)^n}$.

  Chuỗi đều (annuity): $PV = A\cdot\dfrac{1 - (1+r)^{-n}}{r}$, $FV = A\cdot\dfrac{(1+r)^n - 1}{r}$.
\end{congthuc}

\begin{dinhnghia}[title={Phương trình sai phân bậc 1}]
  $y_{t+1} = ay_t + b$. Nghiệm: $y_t = (y_0 - y^*)a^t + y^*$, $y^* = \dfrac{b}{1 - a}$.

  Ổn định: $|a| < 1$.
\end{dinhnghia}

\end{document}
